\documentclass[11pt,a4paper]{article}

% --------------------
% PACKAGES
% --------------------
\usepackage[utf8]{inputenc}
\usepackage[T1]{fontenc}
\usepackage{lmodern}
\usepackage[english]{babel}
\usepackage{geometry}
\usepackage{hyperref}
\usepackage{graphicx}
\usepackage{listings}
\usepackage{xcolor}
\usepackage{fancyhdr}
\usepackage{enumitem}
\usepackage{titlesec}

% --------------------
% PAGE SETTINGS
% --------------------
\geometry{margin=1in}
\pagestyle{fancy}
\fancyhf{}
\fancyhead[L]{\leftmark}
\fancyhead[R]{\thepage}
\setlength{\parskip}{0.5em}

% --------------------
% COLORS AND CODE BLOCKS
% --------------------
\definecolor{codegray}{gray}{1}
\definecolor{codeblue}{rgb}{0.2,0.2,0.6}

\lstset{
	backgroundcolor=\color{codegray},
	basicstyle=\ttfamily\small,
	keywordstyle=\color{codeblue}\bfseries,
	commentstyle=\color{gray}\itshape,
	stringstyle=\color{orange},
	numbers=left,
	numberstyle=\tiny,
	stepnumber=1,
	numbersep=10pt,
	showstringspaces=false,
	breaklines=true,
	frame=single,
	captionpos=b,
	language=C % Change to your preferred language
}

% --------------------
% TITLE
% --------------------
\title{
	\textbf{\Huge The Serotonin Operating System Manual} \\
	\large{C and x86 ASM documentation for interacting with the Serotonin kernel and userspace} \\
	\vspace{0.5cm}
	\vspace{0.3cm}
	Version: 0.1.1 \\
	\vspace{0.5cm}
	\textit{Author: Balazs Monori, Seamus Mullan} \\
	\date{\today}
}
\author{}

\setcounter{secnumdepth}{4}
\titleformat{\paragraph}
{\normalfont\normalsize\bfseries}{\theparagraph}{1em}{}
\titlespacing*{\paragraph}
{0pt}{3.25ex plus 1ex minus .2ex}{1.5ex plus .2ex}


% --------------------
% DOCUMENT
% --------------------
\begin{document}
	
	\maketitle
	\tableofcontents
	\newpage
	
	\section{Preamble}
	This is a manual and is intended to supplement documentation. This manual will not explain every single function available in the operating system, you should also consult documentation for writing programs for Serotonin.
	
	\subsection{Open source acknowledgements}
	Serotonin uses the following open source projects:
	\begin{itemize}
		\item multiboot.h - Copyright (C) 1999,2003,2007,2008,2009,2010  Free Software Foundation, Inc.
		\newline
		Permission is hereby granted, free of charge, to any person obtaining a copy
		of this software and associated documentation files (the "Software"), to
		deal in the Software without restriction, including without limitation the
		rights to use, copy, modify, merge, publish, distribute, sublicense, and/or
		sell copies of the Software, and to permit persons to whom the Software is
		furnished to do so, subject to the following conditions:
		
		The above copyright notice and this permission notice shall be included in
		all copies or substantial portions of the Software.
		
		THE SOFTWARE IS PROVIDED "AS IS", WITHOUT WARRANTY OF ANY KIND, EXPRESS OR
		IMPLIED, INCLUDING BUT NOT LIMITED TO THE WARRANTIES OF MERCHANTABILITY,
		FITNESS FOR A PARTICULAR PURPOSE AND NONINFRINGEMENT.  IN NO EVENT SHALL ANY
		DEVELOPER OR DISTRIBUTOR BE LIABLE FOR ANY CLAIM, DAMAGES OR OTHER LIABILITY,
		WHETHER IN AN ACTION OF CONTRACT, TORT OR OTHERWISE, ARISING FROM, OUT OF OR
		IN CONNECTION WITH THE SOFTWARE OR THE USE OR OTHER DEALINGS IN THE SOFTWARE.
		
		\item GNU Unifont - Copyright (C) GNU Project
		\newline
		GNU GENERAL PUBLIC LICENSE
		Version 2, June 1991
		
		Copyright (C) 1989, 1991 Free Software Foundation, Inc.,
		51 Franklin Street, Fifth Floor, Boston, MA 02110-1301 USA
		Everyone is permitted to copy and distribute verbatim copies
		of this license document, but changing it is not allowed.
		
		Preamble
		
		The licenses for most software are designed to take away your
		freedom to share and change it.  By contrast, the GNU General Public
		License is intended to guarantee your freedom to share and change free
		software--to make sure the software is free for all its users.  This
		General Public License applies to most of the Free Software
		Foundation's software and to any other program whose authors commit to
		using it.  (Some other Free Software Foundation software is covered by
		the GNU Lesser General Public License instead.)  You can apply it to
		your programs, too.
		
		When we speak of free software, we are referring to freedom, not
		price.  Our General Public Licenses are designed to make sure that you
		have the freedom to distribute copies of free software (and charge for
		this service if you wish), that you receive source code or can get it
		if you want it, that you can change the software or use pieces of it
		in new free programs; and that you know you can do these things.
		
		To protect your rights, we need to make restrictions that forbid
		anyone to deny you these rights or to ask you to surrender the rights.
		These restrictions translate to certain responsibilities for you if you
		distribute copies of the software, or if you modify it.
		
		For example, if you distribute copies of such a program, whether
		gratis or for a fee, you must give the recipients all the rights that
		you have.  You must make sure that they, too, receive or can get the
		source code.  And you must show them these terms so they know their
		rights.
		
		We protect your rights with two steps: (1) copyright the software, and
		(2) offer you this license which gives you legal permission to copy,
		distribute and/or modify the software.
		
		Also, for each author's protection and ours, we want to make certain
		that everyone understands that there is no warranty for this free
		software.  If the software is modified by someone else and passed on, we
		want its recipients to know that what they have is not the original, so
		that any problems introduced by others will not reflect on the original
		authors' reputations.
		
		Finally, any free program is threatened constantly by software
		patents.  We wish to avoid the danger that redistributors of a free
		program will individually obtain patent licenses, in effect making the
		program proprietary.  To prevent this, we have made it clear that any
		patent must be licensed for everyone's free use or not licensed at all.
		
		The precise terms and conditions for copying, distribution and
		modification follow.
		
		GNU GENERAL PUBLIC LICENSE
		TERMS AND CONDITIONS FOR COPYING, DISTRIBUTION AND MODIFICATION
		
		0. This License applies to any program or other work which contains
		a notice placed by the copyright holder saying it may be distributed
		under the terms of this General Public License.  The "Program", below,
		refers to any such program or work, and a "work based on the Program"
		means either the Program or any derivative work under copyright law:
		that is to say, a work containing the Program or a portion of it,
		either verbatim or with modifications and/or translated into another
		language.  (Hereinafter, translation is included without limitation in
		the term "modification".)  Each licensee is addressed as "you".
		
		Activities other than copying, distribution and modification are not
		covered by this License; they are outside its scope.  The act of
		running the Program is not restricted, and the output from the Program
		is covered only if its contents constitute a work based on the
		Program (independent of having been made by running the Program).
		Whether that is true depends on what the Program does.
		
		1. You may copy and distribute verbatim copies of the Program's
		source code as you receive it, in any medium, provided that you
		conspicuously and appropriately publish on each copy an appropriate
		copyright notice and disclaimer of warranty; keep intact all the
		notices that refer to this License and to the absence of any warranty;
		and give any other recipients of the Program a copy of this License
		along with the Program.
		
		You may charge a fee for the physical act of transferring a copy, and
		you may at your option offer warranty protection in exchange for a fee.
		
		2. You may modify your copy or copies of the Program or any portion
		of it, thus forming a work based on the Program, and copy and
		distribute such modifications or work under the terms of Section 1
		above, provided that you also meet all of these conditions:
		
		a) You must cause the modified files to carry prominent notices
		stating that you changed the files and the date of any change.
		
		b) You must cause any work that you distribute or publish, that in
		whole or in part contains or is derived from the Program or any
		part thereof, to be licensed as a whole at no charge to all third
		parties under the terms of this License.
		
		c) If the modified program normally reads commands interactively
		when run, you must cause it, when started running for such
		interactive use in the most ordinary way, to print or display an
		announcement including an appropriate copyright notice and a
		notice that there is no warranty (or else, saying that you provide
		a warranty) and that users may redistribute the program under
		these conditions, and telling the user how to view a copy of this
		License.  (Exception: if the Program itself is interactive but
		does not normally print such an announcement, your work based on
		the Program is not required to print an announcement.)
		
		These requirements apply to the modified work as a whole.  If
		identifiable sections of that work are not derived from the Program,
		and can be reasonably considered independent and separate works in
		themselves, then this License, and its terms, do not apply to those
		sections when you distribute them as separate works.  But when you
		distribute the same sections as part of a whole which is a work based
		on the Program, the distribution of the whole must be on the terms of
		this License, whose permissions for other licensees extend to the
		entire whole, and thus to each and every part regardless of who wrote it.
		
		Thus, it is not the intent of this section to claim rights or contest
		your rights to work written entirely by you; rather, the intent is to
		exercise the right to control the distribution of derivative or
		collective works based on the Program.
		
		In addition, mere aggregation of another work not based on the Program
		with the Program (or with a work based on the Program) on a volume of
		a storage or distribution medium does not bring the other work under
		the scope of this License.
		
		3. You may copy and distribute the Program (or a work based on it,
		under Section 2) in object code or executable form under the terms of
		Sections 1 and 2 above provided that you also do one of the following:
		
		a) Accompany it with the complete corresponding machine-readable
		source code, which must be distributed under the terms of Sections
		1 and 2 above on a medium customarily used for software interchange; or,
		
		b) Accompany it with a written offer, valid for at least three
		years, to give any third party, for a charge no more than your
		cost of physically performing source distribution, a complete
		machine-readable copy of the corresponding source code, to be
		distributed under the terms of Sections 1 and 2 above on a medium
		customarily used for software interchange; or,
		
		c) Accompany it with the information you received as to the offer
		to distribute corresponding source code.  (This alternative is
		allowed only for noncommercial distribution and only if you
		received the program in object code or executable form with such
		an offer, in accord with Subsection b above.)
		
		The source code for a work means the preferred form of the work for
		making modifications to it.  For an executable work, complete source
		code means all the source code for all modules it contains, plus any
		associated interface definition files, plus the scripts used to
		control compilation and installation of the executable.  However, as a
		special exception, the source code distributed need not include
		anything that is normally distributed (in either source or binary
		form) with the major components (compiler, kernel, and so on) of the
		operating system on which the executable runs, unless that component
		itself accompanies the executable.
		
		If distribution of executable or object code is made by offering
		access to copy from a designated place, then offering equivalent
		access to copy the source code from the same place counts as
		distribution of the source code, even though third parties are not
		compelled to copy the source along with the object code.
		
		4. You may not copy, modify, sublicense, or distribute the Program
		except as expressly provided under this License.  Any attempt
		otherwise to copy, modify, sublicense or distribute the Program is
		void, and will automatically terminate your rights under this License.
		However, parties who have received copies, or rights, from you under
		this License will not have their licenses terminated so long as such
		parties remain in full compliance.
		
		5. You are not required to accept this License, since you have not
		signed it.  However, nothing else grants you permission to modify or
		distribute the Program or its derivative works.  These actions are
		prohibited by law if you do not accept this License.  Therefore, by
		modifying or distributing the Program (or any work based on the
		Program), you indicate your acceptance of this License to do so, and
		all its terms and conditions for copying, distributing or modifying
		the Program or works based on it.
		
		6. Each time you redistribute the Program (or any work based on the
		Program), the recipient automatically receives a license from the
		original licensor to copy, distribute or modify the Program subject to
		these terms and conditions.  You may not impose any further
		restrictions on the recipients' exercise of the rights granted herein.
		You are not responsible for enforcing compliance by third parties to
		this License.
		
		7. If, as a consequence of a court judgment or allegation of patent
		infringement or for any other reason (not limited to patent issues),
		conditions are imposed on you (whether by court order, agreement or
		otherwise) that contradict the conditions of this License, they do not
		excuse you from the conditions of this License.  If you cannot
		distribute so as to satisfy simultaneously your obligations under this
		License and any other pertinent obligations, then as a consequence you
		may not distribute the Program at all.  For example, if a patent
		license would not permit royalty-free redistribution of the Program by
		all those who receive copies directly or indirectly through you, then
		the only way you could satisfy both it and this License would be to
		refrain entirely from distribution of the Program.
		
		If any portion of this section is held invalid or unenforceable under
		any particular circumstance, the balance of the section is intended to
		apply and the section as a whole is intended to apply in other
		circumstances.
		
		It is not the purpose of this section to induce you to infringe any
		patents or other property right claims or to contest validity of any
		such claims; this section has the sole purpose of protecting the
		integrity of the free software distribution system, which is
		implemented by public license practices.  Many people have made
		generous contributions to the wide range of software distributed
		through that system in reliance on consistent application of that
		system; it is up to the author/donor to decide if he or she is willing
		to distribute software through any other system and a licensee cannot
		impose that choice.
		
		This section is intended to make thoroughly clear what is believed to
		be a consequence of the rest of this License.
		
		8. If the distribution and/or use of the Program is restricted in
		certain countries either by patents or by copyrighted interfaces, the
		original copyright holder who places the Program under this License
		may add an explicit geographical distribution limitation excluding
		those countries, so that distribution is permitted only in or among
		countries not thus excluded.  In such case, this License incorporates
		the limitation as if written in the body of this License.
		
		9. The Free Software Foundation may publish revised and/or new versions
		of the General Public License from time to time.  Such new versions will
		be similar in spirit to the present version, but may differ in detail to
		address new problems or concerns.
		
		Each version is given a distinguishing version number.  If the Program
		specifies a version number of this License which applies to it and "any
		later version", you have the option of following the terms and conditions
		either of that version or of any later version published by the Free
		Software Foundation.  If the Program does not specify a version number of
		this License, you may choose any version ever published by the Free Software
		Foundation.
		
		10. If you wish to incorporate parts of the Program into other free
		programs whose distribution conditions are different, write to the author
		to ask for permission.  For software which is copyrighted by the Free
		Software Foundation, write to the Free Software Foundation; we sometimes
		make exceptions for this.  Our decision will be guided by the two goals
		of preserving the free status of all derivatives of our free software and
		of promoting the sharing and reuse of software generally.
		
		NO WARRANTY
		
		11. BECAUSE THE PROGRAM IS LICENSED FREE OF CHARGE, THERE IS NO WARRANTY
		FOR THE PROGRAM, TO THE EXTENT PERMITTED BY APPLICABLE LAW.  EXCEPT WHEN
		OTHERWISE STATED IN WRITING THE COPYRIGHT HOLDERS AND/OR OTHER PARTIES
		PROVIDE THE PROGRAM "AS IS" WITHOUT WARRANTY OF ANY KIND, EITHER EXPRESSED
		OR IMPLIED, INCLUDING, BUT NOT LIMITED TO, THE IMPLIED WARRANTIES OF
		MERCHANTABILITY AND FITNESS FOR A PARTICULAR PURPOSE.  THE ENTIRE RISK AS
		TO THE QUALITY AND PERFORMANCE OF THE PROGRAM IS WITH YOU.  SHOULD THE
		PROGRAM PROVE DEFECTIVE, YOU ASSUME THE COST OF ALL NECESSARY SERVICING,
		REPAIR OR CORRECTION.
		
		12. IN NO EVENT UNLESS REQUIRED BY APPLICABLE LAW OR AGREED TO IN WRITING
		WILL ANY COPYRIGHT HOLDER, OR ANY OTHER PARTY WHO MAY MODIFY AND/OR
		REDISTRIBUTE THE PROGRAM AS PERMITTED ABOVE, BE LIABLE TO YOU FOR DAMAGES,
		INCLUDING ANY GENERAL, SPECIAL, INCIDENTAL OR CONSEQUENTIAL DAMAGES ARISING
		OUT OF THE USE OR INABILITY TO USE THE PROGRAM (INCLUDING BUT NOT LIMITED
		TO LOSS OF DATA OR DATA BEING RENDERED INACCURATE OR LOSSES SUSTAINED BY
		YOU OR THIRD PARTIES OR A FAILURE OF THE PROGRAM TO OPERATE WITH ANY OTHER
		PROGRAMS), EVEN IF SUCH HOLDER OR OTHER PARTY HAS BEEN ADVISED OF THE
		POSSIBILITY OF SUCH DAMAGES.
		
		END OF TERMS AND CONDITIONS
		
		How to Apply These Terms to Your New Programs
		
		If you develop a new program, and you want it to be of the greatest
		possible use to the public, the best way to achieve this is to make it
		free software which everyone can redistribute and change under these terms.
		
		To do so, attach the following notices to the program.  It is safest
		to attach them to the start of each source file to most effectively
		convey the exclusion of warranty; and each file should have at least
		the "copyright" line and a pointer to where the full notice is found.
		
		<one line to give the program's name and a brief idea of what it does.>
		Copyright (C) <year>  <name of author>
		
		This program is free software; you can redistribute it and/or modify
		it under the terms of the GNU General Public License as published by
		the Free Software Foundation; either version 2 of the License, or
		(at your option) any later version.
		
		This program is distributed in the hope that it will be useful,
		but WITHOUT ANY WARRANTY; without even the implied warranty of
		MERCHANTABILITY or FITNESS FOR A PARTICULAR PURPOSE.  See the
		GNU General Public License for more details.
		
		You should have received a copy of the GNU General Public License along
		with this program; if not, write to the Free Software Foundation, Inc.,
		51 Franklin Street, Fifth Floor, Boston, MA 02110-1301 USA.
		
		Also add information on how to contact you by electronic and paper mail.
		
		If the program is interactive, make it output a short notice like this
		when it starts in an interactive mode:
		
		Gnomovision version 69, Copyright (C) year name of author
		Gnomovision comes with ABSOLUTELY NO WARRANTY; for details type `show w'.
		This is free software, and you are welcome to redistribute it
		under certain conditions; type `show c' for details.
		
		The hypothetical commands `show w' and `show c' should show the appropriate
		parts of the General Public License.  Of course, the commands you use may
		be called something other than `show w' and `show c'; they could even be
		mouse-clicks or menu items--whatever suits your program.
		
		You should also get your employer (if you work as a programmer) or your
		school, if any, to sign a "copyright disclaimer" for the program, if
		necessary.  Here is a sample; alter the names:
		
		Yoyodyne, Inc., hereby disclaims all copyright interest in the program
		`Gnomovision' (which makes passes at compilers) written by James Hacker.
		
		<signature of Ty Coon>, 1 April 1989
		Ty Coon, President of Vice
		
		This General Public License does not permit incorporating your program into
		proprietary programs.  If your program is a subroutine library, you may
		consider it more useful to permit linking proprietary applications with the
		library.  If this is what you want to do, use the GNU Lesser General
		Public License instead of this License.
	\end{itemize}
	
	\section{Hardware Support}
	
	\subsection{Minimum Requirements}
	
	The following specifications are required to run Serotonin:
	\begin{itemize}
		\item x86 CPU with Streaming SIMD Extensions 2 (SSE2) instruction set extension (Intel Pentium 4, AMD Athlon64, VIA C7 or newer)
		\item $\geq$1GiB of RAM
		\item Legacy BIOS with VESA BIOS Extensions 2.0 support
		\item $\geq$TBDGiB of hard disk space
	\end{itemize}
	
	\subsection{Video}
	Serotonin supports the VESA BIOS Extensions, preconfigured for a resolution of 1280x800 with 32-bit true colour (16,777,216 possible colours, excluding the alpha channel).
	
	\subsection{Audio}
	
	\subsubsection{PC Speaker}
	The PC Speaker is capable of generating basic square waves.
	
	\subsubsection{Yamaha YM3812 (OPL2)}
	The YM3812 is a sound integrated circuit, popularised by the AdLib and Sound Blaster sound cards. It is capable of 9 melodic voice polyphony (6 melodic + 5 percussion in rhythm mode). Each voice has 2 operators for Frequency Modulation (FM) and additive synthesis, 4 selectable waveforms per operator (sine, half sine, absolute sine, quarters ine) and per operator 4-bit ADSR (attack, decay, sustain, release).
	
	\subsubsection{Sound Blaster 16}
	TBD
	
	\subsection{Storage}
	
	\subsubsection{IDE}
	Serotonin currently only has support for legacy IDE ATA hard disks.
	
	\subsubsection{Filesystems}
	The following filesystems are supported:
		\begin{itemize}
		\item FAT32
		\item TMPFS/RAMFS
	\end{itemize}
	
	\section{Compilation}
	\subsection{Application Binary Interface}
	Serotonin uses the System V ABI, this includes the calling convention and executable format (ELF). Therefore, it is recommended to use the GNU Compiler Collection (GCC) when compiling both the kernel and and usermode programs.
	
	\subsection{GCC Cross Compilation}
	You should use the i686-elf target for cross compiling from a different operating system other then Serotonin. This will be subject to change when an OS specific toolchain is completed at some point in the future.
	
	\section{The Serotonin Scheduler}
	
	\subsection{Brief Overview}
	The Serotonin Scheduler is a round-robin scheduler for kernel mode (Ring 0, "Supervisor mode") and user mode processes (Ring 3). Kernel mode tasks are cooperative (that is, you must explicity yield the CPU in your code), and usermode tasks are preemptive (the scheduler will preempt your process, but will guarantee at least a quantum, currently $\geq$100ms of CPU time).
	
	\subsection{Process entry point}
	You must provide the entry point of the process. The scheduler will then start execution from the provided pointer, provided that the process is in the scheduler queue, the process is ready and is not blocked. You must also ensure that the loaded binary is in the correct location of memory, that is, kernel mode code should be placed in kernel mapped page area (\texttt{0xC0000000} - \texttt{0xCFFFFFFF}), and the equivalent for user mode (\texttt{0x00400000} - \texttt{0x102FFFFF}).
	
	\subsection{Process stack}
	The scheduler will allocate a stack for you, the size and location depends on the privilege level. In kernel mode processes, a stack of size 64 kibibytes is allocated, located at \texttt{0xF0400000} and growing upwards. In user mode, a stack of size 256 kibibytes is allocated, located at \texttt{0x10400000} and growing upwards.
	
	\subsection{Creating a kernel mode process}
	You can create a new kernel mode process by creating a process and adding it to the scheduler queue:
	\begin{lstlisting}[language=C]
#include "stdio/stdio.h"
#include "schedule/schedule.h"

void my_process(void)
{
	for (int i = 0; i < 5; i++) {
		printf("Hello, kernel world!\n");
		kernel_yield();
	}
	task_exit(EXIT_SIGTERM);
}
		
// kernel_main_high
process_control_block_t *my_pcb = task_create(my_process, "my_process", CPU_KERNEL_MODE);
enqueue(my_pcb);
	\end{lstlisting}
	Note that you have to yield the CPU after you are done with your computation. This is important so that other processes can have CPU time and are not starved. Failing to yield the CPU can result in lockups/freezes, and this is not ideal for the end user experience.
	
	\subsection{Creating a user mode process from kernel mode}
	While in most cases, this is not required, as you should be creating usermode tasks from usermode, it is possible to do so by loading an ELF binary.
	\begin{lstlisting}
# /bin/myusermodetask
.section .rodata
msg: .asciz "Hello, world!\n"

.section .text
.global _start
_start:
	movl $1,   %eax # WRITE
	movl $0,   %ebx # STDOUT
	movl $msg, %ecx # pointer to message
	movl $14,  %edx # length
	int  $0x80      # send syscall to kernel
	
	movl $0,   %eax # EXIT
	movl $15,  %ebx # SIGTERM
	int $0x80       # send syscall to kernel

	\end{lstlisting}
	\begin{lstlisting}[language=C]
#include "stdio/stdio.h"
#include "schedule/schedule.h"
#include "kernel.h"

// kernel_main_high
process_control_block_t *my_user_pcb = kernel_load_elf("/bin/myusermodetask","myusermodetask");
kernel_yield();
	\end{lstlisting}
	The \texttt{kernel\_load\_elf} function will automatically enqueue the task into the scheduler. You can now yield from the current kernel task. Notice how there is no "yield" in the usermode task. This is because all usermode tasks are preempted by the kernel. You can read more about the Serotonin system call interface in the "System Call Interface" chapter.
	
	\subsection{Creating a user mode process from user mode}
	You can create a new user mode process by using the fork and execve system calls. More detail can be viewed in the "System Call Interface" chapter.
	\begin{lstlisting}
.section .text
.global _start
_start:
	movl $4,   %eax # FORK
	int  $0x80      # send syscall to kernel
	
	cmpl $0,   %eax
	je   child      # eax==0 if child
	jg   parent     # eax==child_pid if parent (eax>0)
	
	jmp  error      # eax<0 if fork failed

child:
        # do whatever in child
    
parent:
	# do whatever in parent
	
error:
	# handle error
	
	\end{lstlisting}
	
	\section{System Call Interface}
	\subsection{Brief Overview}
	Serotonin uses the interrupt 0x80 to handle system calls generated by usermode tasks, and arguments are taken from the general purpose registers eax, ebx, ecx and edx. The eax register is used to specify the system call operation, while ebx, ecx and edx are used as arguments for the system call. For example, the exit system call uses eax=0, and ebx is used to specify the signal.

    
	
    \section{Usermode C Library (Not Currently Implemented)}

    Not currently implemented.

    \section{User Interface (Not Currently Implemented)}

    Anything stated is planned, but subject to change.

    \subsection{Taskbar}

    Serotonin's UI contains a simple taskbar with a list of active programs, system information like CPU, GPU, RAM and Storage usage. The taskbar location can be changed in the settings, alongside its width and the colour.

    \subsection{Window Layout}

    \subsubsection{Tiling Window Manager - Ionic}

    Serotonin implements a tiling window manager (WM) called Ionic, named after Ionic bonds. The WM is a simplistic, low fidelity UI made purely for performance. This is advised for lower end hardware nearer to the Hardware Requirements spec. Windows are organised in a grid-like structure and can be moved using \lstinline{MOD + ARROW_KEY}. This WM is enabled by default in the settings.

    \subsubsection{Floating / Stacking Window Manager - Covalence}

    Covalence is a floating WM implemented into Serotonin as an alternative to Ionic. Covalence is much more intuitive for end users familiar with operating systems like MacOS, Windows and other WMs like Gnome or KDE. Named after covalence bonds, this window manager allows users to move, resize, minimize and close windows with their mouse, or optionally with keyboard shortcuts. Naturally this WM is more computationally expensive and must be enabled manually by the user.
    
    \subsection{Running Programs}

    \subsubsection{Quick Run}

    Processes can be opened using the default shortcut \lstinline{MOD + D} which opens the program selection dialog, [insert name here]. The user can then open a previous task using \lstinline{CTRL + NUM}, where \lstinline{NUM} is a number key from 1 to 3.

    \subsubsection{Searching Programs}
    
    The user can also begin to type the name of a program to search the system for it instead. The program is then run by pressing \lstinline{ENTER}. The arrow keys or tab can be used to cycle through the current options, as well as double clicking the program in the list.

    \subsection{Settings}

    Serotonin allows users to customize their desktop experience inside the settings program. Users can change their wallpaper, taskbar colour and location, default window manager and other settings as they get implemented. The menu will render using an Ionic-style window so low end hardware can always run the settings app. This guarantees they can configure their system to perform well.
	
\end{document}
